\documentclass[a4paper]{article}
\usepackage{amsmath}
\usepackage{amsfonts}
\usepackage[affil-it]{authblk}
\usepackage[backend=bibtex,style=numeric]{biblatex}

\usepackage{geometry}
\geometry{margin=1.5cm, vmargin={0pt,1cm}}
\setlength{\topmargin}{-1cm}
\setlength{\paperheight}{29.7cm}
\setlength{\textheight}{25.3cm}
\setlength{\parindent}{0pt}

\begin{document}
% =================================================
\title{Numerical Analysis homework \# 1}

\author{WangHao 3220103613
  \thanks{Electronic address: \texttt{3220103613@zju.edu.cn}}}
\affil{(Information and Computing Sciences, Class 2201), Zhejiang University }


\date{Due time: \today}

\maketitle






% ============================================
\section*{I.}
Consider the bisection method starting with the initial interval $[1.5, 3.5]$. In the following questions �the
interval� refers to the bisection interval whose width
changes across different loops.
\begin{itemize}
    \item What is the width of the interval at the $n$th step?
    \item What is the supremum of the distance between the root $r$ and the midpoint of the interval?
\end{itemize}

\subsection*{Answer of I-a}
Each step will halve the width of the interval. 
Thus, at $n$th step, the width of the interval is $(3.5-1.5) * \left(\frac{1}{2} \right)^n = \left( \frac{1}{2} \right)^{n-1} $.

\subsection*{Answer of I-b}
It is obvious that $\text{left point} \leq r \leq \text{right point}$, 
then the supremum of the distance is $\frac{1}{2} * \text{width of the interval} = \frac{1}{2} * \left( \frac{1}{2} \right)^{n-1} =\left( \frac{1}{2} \right)^{n}$ (note: at $n$th step)

\section*{II.}
In using the bisection algorithm with its initial interval as $[a_0, b_0]$ with $a_0 > 0$, we want to determine the root with its relative error no greater than $\epsilon$. Prove that this goal of accuracy is guaranteed by the following choice of the number of steps, $$n \geq \frac{log(b_0-a_0)-log\epsilon -log(a_0)}{log2} -1$$

\subsection*{Answer of II}
Let the relative error be $E_r$.

Let the root be $\alpha$, we can notice that $0<a<\alpha<b$.

By I, at $n$th step, $ E_r = \frac{(b_0 - a_0)*(\frac{1}{2})^{n+1}}{\alpha}$.

let $E_r\leq\epsilon$, we have $n \geq \frac{log(b_0-a_0)-log\epsilon -log\alpha}{log2} -1$. This mean that when n satisfy the inequality, the relative error is no greater than $\epsilon$.

As $0<a<\alpha$, we have $log(a)<log(\alpha)$. Thus, when $n \geq \frac{log(b_0-a_0)-log\epsilon -log(a_0)}{log2} -1$, we have $n \geq \frac{log(b_0-a_0)-log\epsilon -log\alpha}{log2} -1$, then the relative error is no greater than $\epsilon$.

\section*{III.}
Perform four iterations of Newton's method for the polynomial equation $p(x) = 4x^3-2x^2+3 = 0$ with the starting point $x_0 = -1$. Use a hand calculator and organize results of the iterations in a table.

\subsection*{Answer of III}
$p'(x) = 12x^2-4x$.

Newton's method: $x_{n+1} = x_n - \frac{p(x_n)}{p'(x_n)}$

By iteration, we have:

\renewcommand{\arraystretch}{1.5}
\begin{table}[ht]
  \centering
  \begin{tabular}{|c|c|}
  \hline
  Iteration & Value of $x_n$ \\ \hline
  $x_1$ & $-\frac{13}{16}$    \\ \hline
  $x_2$ & $-\frac{4409}{5720}$  \\ \hline
  $x_3$ & -0.768832384  \\ \hline
  $x_4$ & -0.768828085  \\ \hline
  \end{tabular}
  \caption{results of the iterations}
\end{table}

% ===============================================

\section*{IV.}
Consider a variation of Newton's method in which only the derivative at $x_0$ is used, $$x_{n+1} = x_n - \frac{f(x_n)}{f'(x_0)}$$
Find $C$ and $s$ such that $$e_{n+1} = C(e_n)^s$$
where $e_n$ is the error of Newton's method at step n, s is a constant, and C may depend on $x_n$, the true solution $\alpha$, and the derivative of the function $f$.

\subsection*{Answer of IV}
Taylor's theorem at $x_n$: $$f(x)=f(x_n)+f'(\xi)(x-x_n)$$
where $\xi$ is between $x$ and $x_n$.

Let $x=\alpha$, we have $$0=f(x_n)+f'(\xi)(\alpha-x_n)$$

Substitute it into the variation of Newton's method, we have $$e_{n+1}=e_n-\frac{f'(\xi)}{f'(x_0)}e_n$$.

Thus, we find $C=1-\frac{f'(\xi)}{f'(x_0)}$ and $s=1$, where $\xi$ is between $\alpha$ and $x_n$.

\section*{V.}
Within $(-\frac{\pi}{2},\frac{\pi}{2})$, will the iteration $x_{n+1} =tan^{-1}x_n$ converge?

\subsection*{Answer of V}
Let $f:(-\frac{\pi}{2},\frac{\pi}{2}) \rightarrow (-\frac{\pi}{2},\frac{\pi}{2})$ and $f(x)=tan^{-1}x_n$.

As $f'(x)=\frac{1}{1+x^2} \leq 1$, $f(x)$ is a contraction.

By theorem 1.39, the iteration converges for any choice $x_0 \in (-\frac{\pi}{2},\frac{\pi}{2})$.

\section*{VI.}
Let $p>1$. What is the value of the following continued fraction?
$$x = \frac{1}{p+\frac{1}{p+\frac{1}{p+\dots}}}$$
Prove that the sequence of values converges. (Hint: this can be interpreted as $ x = \lim_{n \to \infty} x_n $, where $x_1 = \frac{1}{p}, x_2 = \frac{1}{p+\frac{1}{p}}, x_3 = \frac{1}{p+\frac{1}{p+\frac{1}{p}}}$, \dots, and so forth. Formulate $x$ as a fixed point of some function.)

\subsection*{Answer of VI}
We interpret the question as $ x = \lim_{n \to \infty} x_n $, where $x_1 = \frac{1}{p}$, and $x_{n+1}=\frac{1}{p+x_n}$.

As $p>1$, it's easy to prove that $0<x_i<1, \forall i \in \mathbb{N}$.

Let $f:(0,1) \rightarrow (0,1)$ and $f(x)=\frac{1}{p+x}$.

As $|f'(x)|=|\frac{1}{(x+p)^2}| \leq 1$, $f(x)$ is a contraction.

By theorem 1.39, the iteration converges for any choice $x_1 \in (0,1)$. Because $x_1$ satisfies the condition, the sequence of values converges.

\section*{VII.}
What happens in problem II if $a_0 < 0 < b_0$? Derive an inequality of the number of steps similar to that in II. In this case, is the relative error still an appropriate measure?

\subsection*{Answer of VII}
Similar to II, we have $n \geq \frac{log(b_0-a_0)-log\epsilon -log|\alpha|}{log2} -1$. This mean that when n satisfy the inequality, the relative error is no greater than $\epsilon$.

As $\alpha$ can be arbitrarily small in this question, $-log|\alpha|$ can also be arbitrarily huge. Extremely, when $\alpha=0$, $log\alpha$ is no longer meaningful. So the relative error is not appropriate. We can use absolute error to estimate the error.

\section*{VIII.}
Consider solving $f(x) = 0 (f \in C^{k+1})$ by Newton's method with the starting point $x_0$ close to a root of multiplicity k. Note that $\alpha$  is a zero of multiplicity $k$ of the function $f$ iff
$$f^{(k)} \neq 0; \forall i<k, f^{(i)}(\alpha)=0$$
\begin{itemize}
    \item How can a multiple zero be detected by examining the behavior of the points $(x_n, f(x_n))$?
    \item Prove that if $r$ is a zero of multiplicity k of the function $f$, then quadratic convergence in Newton's iteration will be restored by making this modification:$$x_{n+1}=x_{n}-k\frac{f(x_n)}{f'(x_n)}$$
\end{itemize}

\subsection*{Answer of VIII-a}
There are 2 ways.

\subsubsection*{Way1}
As $\lim_{x \to \alpha} f'(x) = f'(\alpha) = 0$, 
$\lim_{n \to \infty} (x_{n+1}-x_n) = 0$.


$\lim_{n \to \infty} \frac{f(x_{n+1})-f(x_n)}{x_{n+1}-x_n}=\lim_{n \to \infty}f'(x)=f'(\alpha)=0$.

So we can detect a multiple zero by calculate the slope of $(x_{n+1}, f(x_{n+1}))$ and $(x_n, f(x_n))$.

\subsubsection*{Way2}

Follow the prove of the convergence of Newton's method with the starting point $x_0$ close to a root of multiplicity 1, we can notice that the boundedness of $\frac{f''(\xi)}{f'(x)}$ determines its 2-order.

Now $f'(\alpha) = 0$, $\frac{f''(\xi)}{f'(x)}$ is no longer bounded. We can easily prove that the order of convergence is only 1. 

So we can detect a multiple zero by detect the rate of change.
\subsection*{Answer of VIII-b}
I can't solve that...

\end{document}